\documentclass[11pt]{article}
\usepackage[margin=25mm]{geometry}
\usepackage{fontspec}
\setmainfont{Latin Modern Roman}
\usepackage{amsmath,amssymb,mathtools}
\usepackage{unicode-math}
\setmathfont{Latin Modern Math}
\usepackage{xurl,hyperref}
\hypersetup{colorlinks=true,urlcolor=blue}
\setcounter{secnumdepth}{0}
\setlength{\parindent}{0pt}
\setlength{\parskip}{5pt}
\setlength{\emergencystretch}{3em}
\newcommand{\sref}[2]{\hyperlink{source:#1:#2}{[S#2]}}
\newcommand{\cataloguescope}[1]{\paragraph{Recorded target.}#1}
\newcommand{\R}{\mathbb R}
\newcommand{\C}{\mathbb C}
\newcommand{\N}{\mathbb N}
\newcommand{\Z}{\mathbb Z}
\newcommand{\Q}{\mathbb Q}
\newcommand{\F}{\mathbb F}
\newcommand{\eps}{\varepsilon}
\begin{document}
\section{Affine–Virasoro Derivation of Nekrasov Blow-Up Identities}\label{statement}

\subsection{Definitions and mathematical statement}

\paragraph{Shared notation.}$\mathbb N=\{1,2,\ldots\}$; the standard number fields have their usual meanings. All additional conventions are stated locally.


Write $\Psi(a,m;w,q;\epsilon_1,\epsilon_2)$ for the instanton partition function with surface defect and $Z(a,m;q;\epsilon_1,\epsilon_2)$ for the function without that defect. Here $a$ is the Coulomb parameter, $m$ the four masses, $w$ the defect parameter, $q$ the gauge coupling, and $\epsilon_1,\epsilon_2$ the deformation parameters. The source blow-up relation reads
\[
\Psi(a,m;w,q;\epsilon_1,\epsilon_2)=\sum_{n\in\mathbb Z}
\Psi(a+\epsilon_1n,m;w,q;\epsilon_1,\epsilon_2-\epsilon_1)
Z(a+\epsilon_2n,m;q;\epsilon_1-\epsilon_2,\epsilon_2).
\]
An affine--Virasoro identity decomposes affine Lie-algebra representations into factors involving Virasoro representations.

\paragraph{Catalogue wording retained for traceability.}
Derive Nekrasov's blow-up identity for the $N_f=4$, $\operatorname{SU}(2)$ instanton partition function from the affine–Virasoro vertex-operator-algebra identity under $\frac{\epsilon_1-\epsilon_2}{\epsilon_2}=b^2=-\frac{k+3}{k+2}$, including the required normalization of conformal blocks.
\subsection{Short English statement}
Derive the stated instanton blow-up formula from the affine--Virasoro representation identity, including matching the conformal-block normalizations.
\subsection{Research attempt}
\paragraph{Model, date and investigated scope.}
OpenAI GPT-6 Astra (configured model \texttt{gpt-6-astra}), ultra reasoning; generated 11 October 2026. Identified a coefficientwise divergence under naive unit normalization and checked the deformation/level parameter relation; no affine–Virasoro derivation.

\paragraph{Sources and literature check.}
Read the supplied definition, catalogue wording, source qualifications and reference entries for this question. The historical source/context and inherited status are retained; no fresh literature or novelty certification is claimed. No external theorem is used without proof in this attempt; explicit model assumptions are distinguished from facts about the source model.

\paragraph{A formal normalization obstruction.} Test the displayed blow-up formula under the common but potentially incorrect convention that each instanton-only factor has expansion \(1+O(q)\), independent of the shifted integer \(n\). Then every summand on the right has constant coefficient one. An infinite sum of formal power series is defined coefficientwise only when each coefficient receives finitely many nonzero contributions, or when an additional convergence structure is supplied. Here the putative constant coefficient is
\[
 \sum_{n\in\mathbb Z}1,
\]
so that normalization cannot make the displayed identity meaningful in the ordinary formal power-series ring. Classical factors, lattice-sector prefactors or a different summation convention must be retained. For example exponents \(q^{c_n}\) with \(c_n\to+\infty\) make coefficientwise local finiteness possible, although they do not prove the identity.
\paragraph{Parameter countercheck.} From the stated relation, with \(\epsilon_2\ne0\), \(k\ne-2\), algebra gives
\[
 \frac{\epsilon_1}{\epsilon_2}=-\frac1{k+2}.
\]
Thus finite \(k\) also requires \(\epsilon_1\ne0\). The exceptional \(\epsilon_1=\epsilon_2\) corresponds to \(k=-3\) and \(b^2=0\), where generic conformal-block expressions may require separate limiting treatment.
\paragraph{Remaining derivation.} The first gap is obtaining the actual sector-dependent block normalizations and intertwiners from the affine--Virasoro decomposition, so that the integer sum is well defined and equals the defect partition function. This audit does not declare Nekrasov's identity false; it rejects the stripped normalization. The four masses, defect parameter and all sector coefficients remain to be matched.

\paragraph{Verification and adversarial audit.}
Checked infinite-sum local finiteness and exceptional parameter denominators; treated the failure as a normalization error, not a disproof of the source identity. This is a GPT-6 Astra Ultra self-audit, not independent human verification. Prompt A drafting followed by prompt B scope/dependency/boundary self-audit. Checked infinite-sum local finiteness and exceptional parameter denominators; treated the failure as a normalization error, not a disproof of the source identity. The remaining gap is: Recover the correct integer-sector prefactors, block normalizations and intertwiner coefficients making the blow-up sum meaningful and equal to the source partition function. Unresolved status and zero original-target completion are retained.

Coordinator checked the infinite constant-coefficient obstruction under the stated stripped normalization and the parameter rearrangement. This additional check was model-assisted, not independent human verification.

\paragraph{Tools, process and contributions.}
Python 3 and local file-reading/writing tools for manuscript serialization and structural checks; mathematical calculations were symbolic reasoning, with no numerical experiment or theorem prover. Contribution prompt A was followed by prompt B. GPT-6 Astra with ultra reasoning drafted the research and performed its stated self-audit; any additional coordinator contribution is identified in the verification and audit. A GPT-6-based Codex coordinator assembled the manuscript and checked repository integration; its exact variant and reasoning configuration were not exposed. No human mathematical contribution, independent human proof check or formal proof certification is supplied. Complete derivations and the consulted source references are included here; private reasoning transcripts are not published.

\paragraph{Final status and first remaining gap.}
\textbf{UNRESOLVED}. Recover the correct integer-sector prefactors, block normalizations and intertwiner coefficients making the blow-up sum meaningful and equal to the source partition function.

\section{Completion Estimate}
\noindent\textbf{COMPLETION: 0\%}
The proved calculation is an elementary auxiliary result or a test of a restricted strategy; no new portion of the full source target is established.

\subsection{Sources}
\begin{itemize}
\item[\hypertarget{source:statement:1}{S1}] ULAM.AI and the UnsolvedMath Contributors, UnsolvedMath, catalogue entry. Catalogue curation under CC BY 4.0; definitions and English exposition added locally.
\url{https://huggingface.co/datasets/ulamai/UnsolvedMath}.
\item[\hypertarget{source:statement:2}{S2}] Mini-Workshop: Recent Developments in Representation Theory and Mathematical Physics (2022), Oberwolfach Reports, DOI 10.4171/owr/2022/16, pp. 817--848. Original report and mathematical context.
\url{https://doi.org/10.4171/owr/2022/16}.
\end{itemize}
\label{end:statement}
\end{document}
